% SHARD-ID: LB-11-SOFT-INDEX
% SHARD-TITLE: The soft index: Ward-identity constraints on soft amplitudes
% SHARD-SOURCES: theory/soft-index-r2.md (all sections, incl. the quarantined
%   MERGE PROPOSALS of §7); theory/soft-index-general.md §§0--4;
%   theory/soft-index-g4.md ⟨1⟩1--⟨1⟩7; theory/proto-lsz-match.md;
%   theory/proto-lsz-scalar.md; definitions.md D24, D25;
%   claims/CLAIMS.md rows S-IDX-*, S2-2body-S, D24-VAL, AMP;
%   theory/verdicts/soft-index-r2-critic.md, blitz-proto-lsz-r1.md,
%   d24d3-adjudication-r3/r4/r5.md
% SHARD-CLAIMS: S-IDX-fin-r2, S-IDX-fin-G, S-IDX-G-label,
%   S-IDX-spec-struct-r2, S-IDX-spec-r2, S-IDX-HR-value-r2,
%   S-IDX-MATCH-HS-SEP, S-IDX-D29-value-HS-SEP,
%   S-IDX-PROTO-SCALAR-HS-SEP; proposed D29, D30

% Local semantic macros for this shard (rule 10 permits shard-local macros).
% \statusProposed marks a definition that is NOT yet installed in
% definitions.md: it lives only in a shard's MERGE PROPOSALS block.
\newcommand{\statusProposed}{\statustag{Plum}{proposed, not merged}}
\newcommand{\Aproto}{\mathcal{A}}        % the finite protocol ratio datum
\newcommand{\Rproto}{\mathcal{R}}        % the interacting/free row ratio
\newcommand{\Hsec}{\mathcal{H}}          % a finite weight sector
\newcommand{\rank}{\operatorname{rank}}  % \DeclareMathOperator is preamble-only

\section{The soft index: Ward-identity constraints on soft amplitudes}
\label{sec:soft-index}

\subsection{Why an index, and not a soft factor}

The programme that opened this part of the campaign was the obvious one: look
for a lattice analogue of a universal soft factor --- a single multiplier,
built out of the charges and velocities of the hard legs alone, that
multiplies an $n$-leg amplitude when one extra soft quantum of momentum $k$ is
emitted. That route failed, and it failed for a stated reason: on the lattice
there is no proof that the soft multiplier is independent of the microscopic
source, and there is an explicit local source for which it is not. The
refutation is recorded in \cref{sec:universality}; what survives there is a
criterion and a conditional shape, not a number.

The work in this section replaces that route by a different one. Rather than
asking what the soft emission \emph{is}, we ask what the symmetry algebra
\emph{forces} it to be. The symmetry supplies an exact operator identity in
finite volume --- a Ward identity between a charge, its associated current
zero mode, and the projector onto the range of the charge --- and that identity
constrains a normalised ratio of matrix elements to be exactly one. We call
that ratio the \emph{soft index}. It is an index in the honest sense: it is
an integer (or a one), it is insensitive to the couplings of the Hamiltonian,
to the system size, and to the hard state inside the relevant sector, and its
proof is finite-dimensional linear algebra plus representation theory.

The strategy therefore splits cleanly into two halves, and the split is the
organising idea of this whole section.

\begin{itemize}
\item \textbf{Structure.} What the symmetry alone forces. This half is
  unconditional in finite volume (\cref{thm:soft-index-su2},
  \cref{thm:soft-index-compact-G}, \cref{thm:sector-label}), and it becomes a
  conditional statement about limits of a protocol readout once one asks about
  the soft limit (\cref{thm:soft-index-structural}). Structure never produces
  a number: it produces a shape with an open constant in it.
\item \textbf{Value.} What pins the constant. Values come from on-shell
  two-body data --- the exact spin-$S$ two-magnon phase slope of
  \cref{sec:two-magnon} --- and they enter only after a separately stated
  matching hypothesis that identifies the protocol readout with the on-shell
  multiplier. Where that matching is a hypothesis, the value statement is a
  conditional; where the matching is a theorem, the value statement is a
  theorem.
\end{itemize}

The one substantive advance of the campaign in this area is that on a
restricted but explicitly displayed class of preparations --- the
\emph{separated-preparation} class of \cref{def:separated-preparation} --- the
matching hypothesis is no longer a hypothesis. It is proved, with zero
remainder (\cref{thm:match-separated}), and with it the value follows
unconditionally on that class (\cref{thm:d29-value-separated}) and the
factorisation takes an exact LSZ shape (\cref{thm:proto-scalar}). What remains
open there is not the aggregate identity but its decomposition into named
components; that residual gap is isolated as a single lemma
(\cref{scope:comp-hs}).

\subsection{The protocol datum: a proposed readout, not yet a definition}

Everything in the ``limit'' half of this section is a statement about one
concrete finite-volume number and its subsequential limits. That number, and
the regularity package attached to it, are written down here in full. They are
flagged \statusProposed{} throughout: they live in the merge-proposal block of
the soft-index shard and have deliberately \emph{not} been installed as
numbered definitions in the campaign's definition file, because no adjudicated
round has yet promoted them. Every theorem below that mentions them says so.

\begin{definition}[The windowed, packet-smeared, charge-created ratio readout]
\label{def:protocol-readout}
\statusProposed

Fix a site spin $S$ with $2S\in\N$ and the fully polarised bilinear
ferromagnet $H_S=-J\sum_x(\mathbf S_x\cdot\mathbf S_{x+1}-S^2)$ on an
$N$-site periodic ring, and fix a hard window $I=[a,b]\Subset(0,\pi)$ inside
one physical velocity channel. The readout has six parts.

\emph{(1) Hard register.} Let $\psi_{g,\sigma}$ be a normalised one-magnon
hard packet drawn from $g_\sigma\in C_c^\infty(I)$ on the ring momenta. Write
$m_{\lambda,N}$ for the weight of its sector; on the primitive one-magnon band
$m_{\lambda,N}=NS-1$, while an abstract use carries only
$m_{\lambda,N}/N\to\rho>0$. On the \emph{full} weight sector put
$D_{\lambda,N}=Q_0$, $A_{\lambda,N}=D_{\lambda,N}^\dagger D_{\lambda,N}$ and
$P_{\lambda,N}=D_{\lambda,N}A_{\lambda,N}^{-1}D_{\lambda,N}^\dagger$, and for a
scalar hard band with $v_h\neq0$ record the \emph{measured Ward register}
\begin{equation}
  \ell_{\lambda,N}(h):=
  \frac{\bra{h}Q_0^\dagger P_{\lambda,N}J^-_0\ket{h}}{2iv_h}.
  \label{eq:measured-ward-register}
\end{equation}

\emph{(2) One scale-tied soft packet.} Fix $f\in C_c^\infty((c_1,c_2))$ with
$0<c_1<c_2<1$ and $\norm{f}_2=1$ (or its reflection in the opposite channel),
put $f_\epsilon(k)=\epsilon^{-1/2}f(k/\epsilon)$ and
$\Lambda_N(\epsilon)=(2\pi\Z/N)\cap\supp f_\epsilon$. The \emph{only} soft
insertion is $Q[f_\epsilon]=\sum_{k\in\Lambda_N(\epsilon)}f_\epsilon(k)Q_k$,
and the finite state is
$\Phi_{N,\sigma}(0;\epsilon)=Q[f_\epsilon]\psi_{g,\sigma}$. There is no
independent carrier or width limit and no arbitrary local source.

\emph{(3) Coordinate-kernel register.} A normalised two-excitation vector is
mapped to a symmetric $N\times N$ matrix by sending an occupation coefficient
$c_{xy}$, $x<y$, to $c_{xy}/\sqrt2$ at the entries $(x,y)$ and $(y,x)$ and,
when $S\ge1$, a double-occupancy coefficient $c_{xx}$ to $(x,x)$; the momentum
kernel is the discrete Fourier transform
$\widehat\Phi(k,h)=N^{-1}\sum_{x,y}e^{-i(kx+hy)}\Phi(x,y)$.

\emph{(4) Interacting/free readout.} For a settling time $T$ let
$\widehat\Phi_N(T)$ be the kernel of $e^{-iH_ST}\Phi_{N,\sigma}(0;\epsilon)$
and let
$\widehat\Phi_N^{\rm free}(T)(k,h)=\widehat\Phi_N(0)(k,h)\,
e^{-i[\omega_S(k)+\omega_S(h)]T}$ be the freely propagated reference. With a
hard-column set $K_h\subset I\cap(2\pi\Z/N)$ and an initial ordering matched to
the physical out/in labelling, put
\begin{align}
  d_N(k)&=\sum_{h\in K_h}\abs{\widehat\Phi_N^{\rm free}(T)(k,h)}^2, &
  n_N(k)&=\sum_{h\in K_h}\widehat\Phi_N(T)(k,h)\,
          \overline{\widehat\Phi_N^{\rm free}(T)(k,h)},
  \label{eq:d29-rows}
\end{align}
and, on the domain where $D_N(\epsilon):=\sum_k d_N(k)>0$,
\begin{equation}
  d\mu_{N,\epsilon}(k,h)=
  \frac{\abs{\widehat\Phi_N^{\rm free}(T)(k,h)}^2}{D_N(\epsilon)},
  \qquad
  \Aproto_{N,T,W,\sigma}(\epsilon)=
  \frac{\sum_k n_N(k)}{D_N(\epsilon)}-1,
  \qquad
  \bar k(\epsilon)=\int k\,d\mu_{N,\epsilon}.
  \label{eq:d29-datum}
\end{equation}
The aggregate is a norm projection, not a pointwise division by a hard
amplitude; every finite datum is an explicit finite matrix quantity.

\emph{(5) Exact finite diagnostics.} Record separately the full-sector split
$J^-_k=P_{\lambda,N}J^-_k+(1-P_{\lambda,N})J^-_k$ and both terms of the
current/contact identity, identifying neither with the other; and for an exact
hard vector define the contact defect
$\mathfrak D_N(k,h)=(H_S-\omega_S(k)-\omega_S(h))Q_k\ket{h}_N$, whose time
integral is the exact interacting-minus-free evolution by finite-dimensional
Duhamel. No LSZ exhaustiveness, boundary decay, on-shell matching, value of
the amputation constant, wave operator, or completeness statement is
definitional here.

\emph{(6) Admissible indices and their order.} An index is admissible only if
$N\epsilon(c_2-c_1)>2\pi$, its initial geometry matches the physical channel,
and its settling time lies in a declared settling/recollision sandwich
$[T_{\min},c_{\rm rec}N/v_{\max}]$. Outer limits use admissible sequences
$j\mapsto(N_j,T_j,W_j,\sigma_j)$ at \emph{fixed} $\epsilon$; only after an
outer limit point has been selected is $\epsilon\downarrow0$ taken. At fixed
$N$ the sample set $\Lambda_N(\epsilon)$ is empty once $c_2\epsilon<2\pi/N$,
so the joint sequence is not a limit of this family.
\end{definition}

\provenance{proposed D29 (PROTO)}{%
  \texttt{theory/soft-index-r2.md} §7, step (1.14).(2.1) --- quarantined in
  that shard's merge-proposal block; \emph{not} installed in
  \texttt{definitions.md}}{%
  \texttt{theory/checks/soft\_index\_r2\_check.py} gates SIDXR2-C1--C3, C7--C8
  (finite diagnostics and negative controls)}{%
  \texttt{theory/verdicts/soft-index-r2-critic.md}}

\begin{definition}[The regularity target attached to the readout]
\label{def:protocol-regularity}
\statusProposed

This is a closure property for \cref{def:protocol-readout} and asserts no
existence of an outer limit point. Along a stated admissible outer
subsequence, whenever the scalar datum has an actual cluster family on soft
scales $E\subset(0,\epsilon_0]$ with $0\in\overline E$, require:
\emph{(a) denominator nondegeneracy} --- the aggregate free mass is nonzero,
the normalised row measures are tight on the fixed hard window and the
scale-tied soft support, and rows of zero free mass carry zero measure;
\emph{(b) component compactness} --- the row measures, $m_{\lambda,N}/N$,
$\ell_{\lambda,N}$, and every component named by any separately assumed
decomposition have a further cluster point, with limits written
$\mu_{*,\epsilon}$, $\rho>0$, $\ell_h$, and in particular
$c_1\epsilon\le\abs{\bar k_*(\epsilon)}\le c_2\epsilon$ in a one-sided
channel; \emph{(c) regularity only} --- each cluster family admits the
continuous or $C^1$ extension to $\epsilon=0$ required at its point of use,
with locally uniform bounds and equicontinuous first difference quotients;
\emph{(d) no hidden matching} --- there is no equality with a chamber
coefficient, no on-shell eigenvector, no physical multiplier, no LSZ
exhaustiveness, no value or reality condition for the amputation constant, and
no formula for $\ell_h$ in terms of a composite charge; \emph{(e) optional
convenience} --- full-family convergence may be assumed, buying uniqueness of
the selected limit function only.
\end{definition}

\begin{scope}[Regularity cannot supply a value]
\label{scope:regularity-no-value}
Clause (e) is stated with its own limitation because the limitation is the
point. For any regular family with $1+\Aproto(\epsilon)\neq0$ the
transformation
\begin{equation}
  1+\Aproto(\epsilon)\ \longmapsto\ [1+\Aproto(\epsilon)]\,e^{ic\epsilon}
  \label{eq:mobius-shift}
\end{equation}
preserves every clause (a)--(d) and shifts the phase jet
$\partial_\epsilon\arg(1+\Aproto)|_0$ by the arbitrary real constant $c$. Any
argument that reads a numerical coefficient out of the regularity package
alone is therefore invalid: a separate value supplier must appear at the exact
leaf where the coefficient is fixed. In this section that supplier is always
the exact two-magnon phase slope of \cref{sec:two-magnon}, and never a
regularity clause.
\end{scope}

\provenance{proposed D30 (TGT)}{%
  \texttt{theory/soft-index-r2.md} §7, step (1.14).(2.2); the value-shift
  obstruction is step (1.7)}{%
  \texttt{theory/checks/soft\_index\_r2\_check.py} gate SIDXR2-C6
  (\textsc{tgt-mobius})}{%
  \texttt{theory/verdicts/soft-index-r2-critic.md}}

\subsection{The finite soft-index theorem on the \texorpdfstring{$SU(2)$}{SU(2)} ring}

The first theorem is the anchor of the whole section: it is exact, finite, and
free of every limiting hypothesis. Its content is that the charge, its current
zero mode, and the projector onto the range of the charge are locked together
by the algebra, and that the resulting normalised ratio is one.

The one thing the reader must carry away is that the identity comes in
\emph{two registers} which look similar and are not interchangeable. In the
first, the Gram operator $A=D^\dagger D$ is formed on the \emph{whole} weight
sector and inverted there. In the second, the charge is restricted to the
highest-weight subspace first, and only then is the Gram operator formed; it
is then a scalar. Substituting the scalar of the second register into the
first is false as soon as the sector contains two or more magnons --- a mistake
the campaign made once and now guards against with a dedicated red mode.

\begin{theorem}[Finite soft index on an $SU(2)$ ring]
\label{thm:soft-index-su2}
\statusProved

Let a finite periodic ring carry an on-site $SU(2)$ action and an
$SU(2)$-invariant finite-range Hamiltonian. Let $Q_0=S^-$, let
$(J^+_0,J^z_0,J^-_0)$ be its vector-current zero modes, and let
$\Hsec_{\lambda,N}$ be the \emph{full} weight sector on which
$S^z=m_{\lambda,N}>0$. Put
\[
  D_{\lambda,N}:=Q_0|_{\Hsec_{\lambda,N}},\qquad
  A_{\lambda,N}:=D_{\lambda,N}^\dagger D_{\lambda,N}
  \quad\text{on all of }\Hsec_{\lambda,N},
\]
let $P_{\lambda,N}$ project onto $\ran D_{\lambda,N}$, and let $\Pi_{\rm hw}$
project $\Hsec_{\lambda,N}$ onto
$K_{\lambda,N}:=\ker S^+\cap\Hsec_{\lambda,N}$. Then $A_{\lambda,N}$ is
strictly positive, and for every $\psi\in K_{\lambda,N}$:

\emph{(i) Full-sector Gram-inverse register.}
\begin{equation}
  P_{\lambda,N}J^-_0\psi=2\,D_{\lambda,N}A_{\lambda,N}^{-1}J^z_0\psi.
  \label{eq:fin1}
\end{equation}

\emph{(ii) Separately restricted highest-weight register.} With $D_{\rm hw}$
the restriction of $D_{\lambda,N}$ to $K_{\lambda,N}$ and $P_{\rm hw}$ the
projector onto its range,
\begin{equation}
  D_{\rm hw}^\dagger D_{\rm hw}=2m_{\lambda,N}\mathbb 1,
  \qquad
  P_{\rm hw}J^-_0\psi=\frac{1}{m_{\lambda,N}}\,Q_0\,\Pi_{\rm hw}J^z_0\psi.
  \label{eq:fin2}
\end{equation}

\emph{(iii) Primitive one-magnon residue.} For a nonzero-momentum primitive
one-magnon vector $\ket{h}_N$ in the spin-$S$ ferromagnetic band, where
$m_{\lambda,N}=NS-1$, the projector in \cref{eq:fin2} may be removed and
\begin{equation}
  J^z_0\ket{h}_N=i\,v_S(h)\ket{h}_N,
  \qquad
  \bra{h}Q_0^\dagger P_{\lambda,N}J^-_0\ket{h}_N=2i\,v_S(h).
  \label{eq:fin3}
\end{equation}
Equivalently, the normalised Ward register
\cref{eq:measured-ward-register} equals $\ell_{\lambda,N}(h)=1$.

\emph{(iv) Pure charge-created anchor.} If $\psi_g$ is a one-magnon packet of
exact band vectors and $\Phi_0:=Q_0\psi_g$ is the \emph{pure} zero-mode-created
state, then its interacting/free row ratio in the readout of
\cref{def:protocol-readout} at soft row $k=0$, with nonzero free denominator,
is exactly one for every finite $N,T,W,\sigma$ and every site spin $S$:
\begin{equation}
  \Rproto_{\Phi_0}(0)=1.
  \label{eq:fin8}
\end{equation}
\end{theorem}

\begin{scope}[What the two registers do and do not say]
\label{scope:soft-index-su2}
The two registers are not interchangeable, and the difference is not
cosmetic. In \cref{eq:fin1} the inverse Gram operator stays on the full source
sector, because $J^z_0\psi$ need not remain highest weight once the sector
contains two or more magnons; the scalar $2m_{\lambda,N}$ of \cref{eq:fin2}
may not be pulled through $A_{\lambda,N}^{-1}$. In particular the string
$m_{\lambda,N}^{-1}Q_0J^z_0$, without the highest-weight projector, is
\emph{false} for two or more magnons, where the cross-register error is
$O(1)$ while \cref{eq:fin1} is exact; that string appears nowhere as a claim.

The anchor \cref{eq:fin8} is an anchor for the object just named and nothing
else. It does not assert that the adjoined $k=0$ row of the running family
$Q[f_\epsilon]\psi_g$ equals one at $S\ge1$ --- a direct computation gives
$2.1633+0.1527i$ for that row at $S=1$, $N=7$, against one to
$9.5\times10^{-16}$ for the pure charge-created row --- and it implies nothing
about $\epsilon\downarrow0$, since the admissible family contains no
fixed-$N$ continuous path to that row. At $S=1/2$ the same row identity holds
for arbitrary two-magnon input, so there it is a readout tautology rather than
independent evidence for a soft zero.

The theorem makes no assumption of the protocol readout, its regularity
package, wave operators, completeness, integrability, or any soft limit.
\end{scope}

\paragraph{Idea of proof.} Three short steps. First, positivity: decomposing
the finite $SU(2)$ representation into spin-$j$ irreducibles, a weight-$m$
vector satisfies $S^+S^-=(j+m)(j-m+1)\ge 2m>0$, so $A_{\lambda,N}$ is
invertible on the entire sector and the polar formula
$P_{\lambda,N}=D_{\lambda,N}A_{\lambda,N}^{-1}D_{\lambda,N}^\dagger$ holds.
Second, the Ward step: covariance of the invariant finite-range current gives
$[S^+,J^-_0]=2J^z_0$, so on a highest-weight $\psi$,
$D_{\lambda,N}^\dagger J^-_0\psi=S^+J^-_0\psi=2J^z_0\psi$; substituting this
into the polar formula is exactly \cref{eq:fin1}. Restricting the domain
first inserts the source projector $\Pi_{\rm hw}$ into the adjoint and turns
the Gram operator into the scalar $2m_{\lambda,N}$, giving \cref{eq:fin2}.
Third, the residue: for $h\neq0$ the raised current $J^+_0\ket{h}_N$ vanishes
because the zero-mode current preserves momentum while its raised target is
the momentum-zero vacuum, so $J^z_0\ket{h}_N$ is itself highest weight and the
projector drops; the continuity equation
$-[\omega_S(h+k)-\omega_S(h)]=(e^{ik}-1)\bra{h+k}J^z_k\ket{h}_N$ has removable
value $i\partial_h\omega_S(h)=iv_S(h)$ at $k=0$, and pairing with
$Q_0\ket{h}_N$, whose squared norm is $2m_{\lambda,N}$, cancels the weight and
leaves $2iv_S(h)$. The anchor \cref{eq:fin8} is immediate from
$[H,Q_0]=0$ and $\omega_S(0)=0$: the interacting and free kernels coincide row
by row.

\provenance{S-IDX-fin-r2}{%
  \texttt{theory/soft-index-r2.md} steps (1.3)--(1.5); positivity (1.3).(2.1),
  full register (1.3).(2.4), restricted register (1.3).(2.5), residue
  (1.3).(2.7)--(2.8), anchor (1.4)}{%
  \texttt{theory/checks/soft\_index\_r2\_check.py} gates SIDXR2-C1--C3, green
  exit $0$ with eight public red modes; the \texttt{-{}-red-register-trap}
  mode executes the standing cross-register mutation and propagates its
  outcome}{%
  \texttt{theory/verdicts/soft-index-r2-critic.md} (\textsc{promotable now};
  proof independently recomputed by the critic, promoted after the mandated
  checker repair)}

\subsection{The same theorem for every compact group}

Nothing in the previous proof used $SU(2)$ beyond one root ladder. Written in
that language the theorem generalises verbatim to any compact Lie group, on
any finite graph or periodic lattice, one represented root at a time. This is
worth stating in full because it makes visible exactly which structural
ingredients are load-bearing --- a root triple, an adjoint-covariant current,
and an occupied weight sector --- and therefore exactly where the statement has
no content at all.

\begin{theorem}[Root-wise finite soft index for a compact group]
\label{thm:soft-index-compact-G}
\statusProved

Let a finite graph or a finite periodic lattice carry a finite-dimensional
on-site unitary representation $U_N=u^{\otimes N}$ of a compact Lie group $G$
and a $G$-invariant finite-range Hamiltonian, and put
$Q(X)=\sum_x q_x(X)$ for $X\in\mathfrak g_\C$. Suppose given a displayed
complex-linear current zero mode
$J_0:\mathfrak g_\C\to\mathrm{End}(\Hilb_N)$ satisfying the adjoint covariance
\begin{equation}
  [Q(X),J_0(Y)]=J_0([X,Y]);
  \label{eq:adjoint-covariance}
\end{equation}
on a periodic lattice $J_0$ may be taken to be the displayed directional
zero-momentum flux. Let $\alpha$ be a represented root of the semisimple Lie
algebra of the identity component $G^0$, with a compact-compatible triple
$[E_\alpha,F_\alpha]=H_\alpha$, $[H_\alpha,E_\alpha]=2E_\alpha$,
$[H_\alpha,F_\alpha]=-2F_\alpha$, $Q(E_\alpha)^\dagger=Q(F_\alpha)$ and
$Q(H_\alpha)$ Hermitian. For an occupied positive coroot weight $\lambda$ put
\[
  \Hsec^\alpha_{\lambda,N}:=\ker(Q(H_\alpha)-\lambda),\qquad
  D_\alpha:=Q(F_\alpha)|_{\Hsec^\alpha_{\lambda,N}},\qquad
  A_\alpha:=D_\alpha^\dagger D_\alpha,
\]
let $P_\alpha$ project onto $\ran D_\alpha$, let
$K^\alpha_{\lambda,N}:=\ker Q(E_\alpha)\cap\Hsec^\alpha_{\lambda,N}$ with
projector $\Pi^\alpha_{\rm hw}$, and let $D_{\rm hw}$, $P_{\rm hw}$ be the
correspondingly restricted objects. Then $A_\alpha>0$, and for every
$\psi\in K^\alpha_{\lambda,N}$:
\begin{align}
  P_\alpha J_0(F_\alpha)\psi
    &=D_\alpha A_\alpha^{-1}J_0(H_\alpha)\psi,
    \label{eq:hg-full}\\[2pt]
  D_{\rm hw}^\dagger D_{\rm hw}=\lambda\mathbb 1,
  \qquad
  P_{\rm hw}J_0(F_\alpha)\psi
    &=\frac{1}{\lambda}\,Q(F_\alpha)\,\Pi^\alpha_{\rm hw}J_0(H_\alpha)\psi.
    \label{eq:hg-restricted}
\end{align}
Whenever $\braket{\psi,J_0(H_\alpha)\psi}$ is nonzero, the normalised Ward
index is exactly one:
\begin{equation}
  \frac{\braket{D_\alpha\psi,\;P_\alpha J_0(F_\alpha)\psi}}
       {\braket{\psi,\;J_0(H_\alpha)\psi}}=1.
  \label{eq:hg-index}
\end{equation}
If a possibly disconnected compact subgroup $K\le G$ is used for a Schur or
sector statement, the required global covariance
$U_N(k)J_0(Y)U_N(k)^\dagger=J_0(\Ad_kY)$, $k\in K$, is displayed as a
hypothesis; infinitesimal covariance alone is not used for disconnected
elements.
\end{theorem}

\begin{scope}[Three fences: where this theorem has no rows]
\label{scope:soft-index-compact-G}
The theorem has rows only for \emph{represented roots} of the semisimple Lie
algebra of $G^0$, and this is a genuine restriction with three separate
consequences. A central-torus direction has no root ladder, hence no row at
all. A finite group has no nonzero Lie-algebra current component, hence no
Lie-current row. And an arbitrary finite graph supplies no canonical
direction and no zero mode: on a periodic lattice one may use the displayed
directional flux, but on a general graph the current map is an input, not
something the geometry provides. A chosen finite-ring periodisation of a
termwise invariant interaction is an \emph{example} supplying
\cref{eq:adjoint-covariance}; it is not part of the abstract hypotheses, and
no canonical half-line cut on a circle is asserted.

As in the $SU(2)$ case, the full-sector Gram-inverse register
\cref{eq:hg-full} and the root-highest-restricted scalar register
\cref{eq:hg-restricted} are not interchangeable: $J_0(H_\alpha)\psi$ need not
remain root-highest, so $\lambda$ cannot replace $A_\alpha$ in
\cref{eq:hg-full}.
\end{scope}

\paragraph{Idea of proof.} Restrict the finite unitary representation to the
root $\mathfrak{sl}_2$ and use complete reducibility,
$\Hilb_N\cong\bigoplus_{n\ge0}V_n\otimes M_{n,\alpha}$. On the weight-$\lambda$
line of $V_n$,
\[
  Q(E_\alpha)Q(F_\alpha)=\tfrac14(n+\lambda)(n-\lambda+2)\,\mathbb 1
  \ \ge\ \lambda>0,
\]
so $A_\alpha$ is strictly positive on the \emph{entire} source sector and
$P_\alpha=D_\alpha A_\alpha^{-1}D_\alpha^\dagger$. Adjoint covariance gives
$D_\alpha^\dagger J_0(F_\alpha)\psi=Q(E_\alpha)J_0(F_\alpha)\psi
=J_0(H_\alpha)\psi$ for root-highest $\psi$, and substituting into the polar
formula yields \cref{eq:hg-full}. For the restricted register,
$Q(E_\alpha)Q(F_\alpha)\phi=(Q(F_\alpha)Q(E_\alpha)+Q(H_\alpha))\phi
=\lambda\phi$ on $K^\alpha_{\lambda,N}$, and the adjoint of a
restricted-domain map carries the source projector. The index
\cref{eq:hg-index} needs no further work at all: since $D_\alpha\psi$ lies in
$\ran D_\alpha$, the projector is invisible in the pairing and
$\braket{D_\alpha\psi,P_\alpha J_0(F_\alpha)\psi}
=\braket{\psi,D_\alpha^\dagger J_0(F_\alpha)\psi}
=\braket{\psi,J_0(H_\alpha)\psi}$; divide. No injectivity, rank-one property,
or scalarity of a reduced operator is used beyond the positivity already
proved.

\provenance{S-IDX-fin-G}{%
  \texttt{theory/soft-index-general.md} §§1--2 (statement §1.2, positivity and
  full register, restricted register, normalised index);
  \texttt{theory/soft-index-2d.md} §§1--2}{%
  \texttt{theory/checks/soft\_index\_2d\_check.py} gates SIDX2D-C0--C3 and
  \texttt{-{}-red-scalar-full}; the general-$G$ and $SU(2)$ reference gates;
  fences computed as G5-C1-\textsc{abelian}, G5-C2-\textsc{finite},
  G5-C3-\textsc{global-form}}{%
  \texttt{theory/verdicts/soft-index-r2-critic.md} (reduction to
  \cref{thm:soft-index-su2} verified)}

\subsection{The soft index as a difference of superselection labels}

There is a second, complementary statement in the general compact setting,
and it is the one that deserves the name ``index'' most literally. Instead of
asking about operator matrix elements it asks about \emph{labels}: the
charge-created and current-created vectors sit in a central character sector,
and the soft datum is the difference between the target label and the source
label. That difference is a character, it takes integer values on
cocharacters, and it does not depend on the ring size, the hard vector, the
Hamiltonian couplings, or any inventory of scattering channels.

\begin{theorem}[Sector-label integrality of the soft index]
\label{thm:sector-label}
\statusProved

Let a finite periodic ring carry a compact, possibly disconnected, on-site Lie
symmetry $G$, let $\alpha$ be a vacuum label with unbroken subgroup
$H_\alpha$, write $Z_\alpha:=Z(H_\alpha)$, and define the \emph{effective
unbroken centre}
\begin{equation}
  K_u:=\{z\in Z_\alpha: u(z)\in U(1)\mathbb 1\},
  \qquad
  Z_{\alpha,\rm eff}:=Z_\alpha/K_u,
  \qquad
  \Lambda_\alpha^*:=\Hom(Z_{\alpha,\rm eff},U(1)).
  \label{eq:effective-centre}
\end{equation}
All occupied $Z_\alpha$ characters have the same restriction to $K_u$, so they
form an affine torsor whose differences lie in $\Lambda_\alpha^*$; this
character group carries the free lattice of the connected central torus and
may in addition carry a finite torsion part contributed by disconnected
components. Let a nonzero represented Lie component $X\in\mathfrak g_\C$ span
a $Z_\alpha$-character line of character $\chi_X\in\Lambda_\alpha^*$, put
$Q_X=Q(X)$, let $J_X$ be the displayed finite-ring current component, and
assume the displayed global covariance
\begin{equation}
  U_N(z)Q_XU_N(z)^\dagger=\chi_X(z)Q_X,
  \qquad
  U_N(z)J_XU_N(z)^\dagger=\chi_X(z)J_X,
  \qquad z\in Z_\alpha.
  \label{eq:global-covariance}
\end{equation}
Then for every occupied central character $\gamma$, with $\Pi_\gamma$ its
projector,
\begin{equation}
  \Pi_\delta Q_X\Pi_\gamma=\delta_{\delta,\gamma\chi_X}\,Q_X\Pi_\gamma,
  \qquad
  \Pi_\delta J_X\Pi_\gamma=\delta_{\delta,\gamma\chi_X}\,J_X\Pi_\gamma,
  \label{eq:label-blocks}
\end{equation}
so that the finite sector-label index is the character difference
\begin{equation}
  \mathrm{ind}_{Z_\alpha}(X;\gamma):=(\gamma\chi_X)\gamma^{-1}
  =\chi_X\in\Lambda_\alpha^*,
  \label{eq:label-index}
\end{equation}
and for every cocharacter $\eta:U(1)\to Z_{\alpha,\rm eff}^0$,
\begin{equation}
  \chi_X(\eta(e^{i\theta}))=e^{in_\eta(\chi_X)\theta},
  \qquad
  \Delta q_\eta=n_\eta(\chi_X)\in\Z.
  \label{eq:label-integrality}
\end{equation}
A disconnected effective centre may additionally contribute the finite torsion
label retained in \cref{eq:label-index}; it does not create an integer circle
coordinate. Finally, for an exact hard eigenvector $\psi_h$ with
$Q_X\psi_h\neq0$, the pure charge-created readout of \cref{eq:fin8} is one
whenever its denominator is nonzero.
\end{theorem}

\begin{scope}[What the label theorem does \emph{not} imply]
\label{scope:sector-label}
The most important negative statement here is that
\cref{eq:label-blocks} does \emph{not} imply the projected-current operator
identity. Writing $D_\gamma:=Q_X\Pi_\gamma$ and $P_{D,\gamma}$ for the
projector onto its range, the theorem gives only the inclusion
\begin{equation}
  \ran D_\gamma\subseteq\ran\Pi_{\gamma\chi_X},
  \qquad
  P_{D,\gamma}\le\Pi_{\gamma\chi_X},
  \label{eq:descendant-finer}
\end{equation}
and the inclusion is strict in an explicit instance. In the four-qutrit
fundamental $SU(3)$ ferromagnet one finds
$\dim\Hsec_{5/3}=8$, $\dim\Hsec_{2/3}=24$, $\rank D=8$, and the live
replacement defect
$\norm{(\Pi_{2/3}-P_D)J_X\ket{h;2}}=\sqrt{8/3}$: replacing the
descendant-range projector by the central-sector projector changes the current
vector. The underdetermination behind this is abstract, not accidental: for a
one-dimensional source line and a two-dimensional target of the same
character, $D\psi=e_1$ and $J\psi=ae_1+be_2$ have the same source and target
labels for every $a,b$, while $P_DJ\psi=ae_1$ varies. Sector labels alone do
not determine an operator coefficient.

Two further degeneracies. If $X$ is central its adjoint character is trivial
and the label \cref{eq:label-index} is zero --- which is still distinct from a
root row, since a central direction has no ladder. If $G$ is finite there is
no nonzero $X$ and the theorem has no Lie-current instance at all; an
independently supplied finite torsion label, or a finite-string endpoint
identity, is a different statement.
\end{scope}

\paragraph{Idea of proof.} If $\psi=\Pi_\gamma\psi$ then
\cref{eq:global-covariance} gives
$U_N(z)Q_X\psi=\gamma(z)\chi_X(z)Q_X\psi$ and likewise for $J_X$, so the
image lies entirely in the $\gamma\chi_X$ block; Haar orthogonality of
distinct characters --- equivalently
$\Pi_\delta=\int_{Z_\alpha}\overline{\delta(z)}\,U_N(z)\,dz$ --- kills every
other block and proves \cref{eq:label-blocks}. Dividing the target label by
the source label proves \cref{eq:label-index}, the common affine restriction
on $K_u$ cancelling; and every continuous character of $U(1)$ is
$e^{i\theta}\mapsto e^{in\theta}$ for a unique integer $n$, which is
\cref{eq:label-integrality}. Since $H_N$ commutes with the central action it
preserves each $\Pi_\gamma$, so the label difference cannot change under
finite-volume dynamics. This is compact representation covariance, not a
statement about quantisation of an expectation value.

\provenance{S-IDX-G-label}{%
  \texttt{theory/soft-index-g4.md} steps (1.1)--(1.7); headline statement at
  (1.3), the pure zero-mode anchor at (1.4), the exact route boundary at
  (1.5), the $SU(2)$ reduction at (1.6), the four-site $SU(3)$ instance at
  (1.7); restated in \texttt{theory/soft-index-general.md} §4}{%
  \texttt{theory/checks/soft\_index\_g4\_check.py} gates G4-C0--C4;
  \texttt{theory/checks/soft\_index\_g\_boundary\_check.py} gates G5-C1--C3;
  \texttt{theory/checks/soft\_index\_r2\_check.py} gates SIDXR2-C2--C3}{%
  \texttt{theory/verdicts/soft-index-r2-critic.md}}

\subsection{The structural limit law, and the one remaining gap in one dimension}

We now leave exact finite volume. The question is what the symmetry forces
about \emph{limits} of the readout of \cref{def:protocol-readout} as the ring,
the settling time and the window are taken out and the soft scale $\epsilon$
is then sent to zero. The honest register for this is the \emph{actual limit
point}: a common outer subsequence along which, for every $\epsilon$ in a set
$E\downarrow0$, the scalar datum, the row measures, the hard Ward registers,
and every component named below all converge in the topologies required. The
phrase asserts nothing about existence of such a subsequence. Along it the
ordered soft variable is the packet mean $\bar k_*(\epsilon)$, not the carrier
of an unsmeared plane wave, and it obeys
$c_1\epsilon\le\abs{\bar k_*(\epsilon)}\le c_2\epsilon$, with sign fixed by
the chosen one-sided profile.

The whole content of the structural theorem rests on one hypothesis, and that
hypothesis is the single remaining gap of the one-dimensional programme. It is
displayed here as a named object so that it can never be smuggled into a
definition.

\begin{definition}[Hypothesis \textsc{(proto-lsz)}: an exhaustive component decomposition]
\label{def:proto-lsz}
In the packet and amputation normalisation of \cref{def:amputation-convention}
(\cref{sec:universality}), after the outer limit and uniformly for $h\in I$,
the limit datum decomposes exhaustively as
\begin{equation}
  \Aproto_*=\Aproto_*^{\rm desc}+\Aproto_*^{\perp}
            +\Aproto_*^{\rm dir}+\Aproto_*^{\partial},
  \label{eq:proto-lsz-split}
\end{equation}
naming respectively the descendant, orthogonal-current, direct-contact, and
two window-boundary-gradient contributions, where the descendant term has the
external-flux form
\begin{equation}
  \Aproto_*^{\rm desc}(\epsilon)=
  \int(e^{ik}-1)\,L(k,h)\,[\,2iv_h\ell_h\,]\;d\mu_{*,\epsilon}(k,h),
  \label{eq:proto-lsz-desc}
\end{equation}
with $L$ uniformly $C^1$ and carrying exactly the $h$-profile
$L(0,h)=\aleg(\rho)\,(-i\sgn(v_h-v_s)/v_h)$ of
\cref{def:no-contact-class}(3b), and
\begin{equation}
  \Aproto_*^{\perp}=O(\epsilon^2),\qquad
  \Aproto_*^{\rm dir}=O(\epsilon^2),\qquad
  \Aproto_*^{\partial}=o(\epsilon).
  \label{eq:proto-lsz-bounds}
\end{equation}
The hypothesis also carries the antecedents that the relevant class
$\classSW(\rho)$ is nonempty and that the asymptotic one-magnon kernel exists
--- solely so that $\aleg(\rho)$ is defined --- and that the measured hard Ward
register converges to a finite $\ell_h$, with $\ell_h=1$ only for the
primitive one-magnon band proved in \cref{thm:soft-index-su2}.
\end{definition}

This hypothesis is uninstantiated. It carries the exhaustiveness, the
kinematic normalisation, the reduced-channel regularity, the no-contact
condition, and the finite-window boundary limit as \emph{explicit assumptions};
it does not make exhaustiveness true by definition, and no current claim
proves it for the adjudicated fixed-time readout.

\begin{theorem}[Structural law for actual limit points]
\label{thm:soft-index-structural}
\statusSketch

Assume an actual limit point of \cref{def:protocol-readout} in the sense
above, the regularity package \cref{def:protocol-regularity}, and hypothesis
\textsc{(proto-lsz)} of \cref{def:proto-lsz}. Then every such limit point
obeys
\begin{equation}
  \Aproto_*(\epsilon)=
  2i\,\aleg(\rho)\,\sgn(v_h-v_s)\,\ell_h\,\bar k_*(\epsilon)+o(\epsilon),
  \label{eq:spec-structural}
\end{equation}
and hence has the continuous Adler extension $\Aproto_*(0)=0$. Before any
on-shell matching its phase jet is only
\begin{equation}
  \frac{\arg(1+\Aproto_*(\epsilon))}{\bar k_*(\epsilon)}
  \longrightarrow
  2\,\mathrm{Re}\,\aleg(\rho)\,\sgn(v_h-v_s)\,\ell_h ,
  \label{eq:spec-phase-jet}
\end{equation}
which is not a numerical law.
\end{theorem}

\begin{scope}[Exactly what is proved and what is missing]
\label{scope:soft-index-structural}
What is proved is a complete conditional implication: given the displayed
hypotheses, the algebra of \cref{eq:spec-structural} is exact. What is missing
is any instance. Hypothesis \textsc{(proto-lsz)} is not proved for the
adjudicated fixed-time readout; in particular the two window-boundary
gradients and the identification of the fixed-time charge insertion with a
constructed scattering channel are not controlled by the wave-operator results
of \cref{sec:wave-operators}, which explicitly exclude that interface. The
soft-leg amputation constant $\aleg(\rho)$ appearing in
\cref{eq:spec-structural} is the \emph{open} class constant of
\cref{def:no-contact-class}(3b): no value, and in particular no $1/\rho$ law,
is claimed here at any density including $\rho=1/2$. The regularity package
cannot repair this by \cref{scope:regularity-no-value}. Finally, existence of
an actual limit point is not asserted.
\end{scope}

\paragraph{What the argument does.} Uniformly on the packet support,
$e^{ik}-1=ik+O(\epsilon^2)$ and
$L(k,h)=\aleg(\rho)(-i\sgn(v_h-v_s)/v_h)+O(\epsilon)$; multiplying the three
leading factors in \cref{eq:proto-lsz-desc} gives
$(ik)\cdot\aleg(-i\sgn/v_h)\cdot 2iv_h\ell_h
=2i\aleg\sgn\,\ell_h\,k$, and integration replaces $k$ by the packet mean.
Every cross term is $O(\epsilon^2)$ uniformly because $I\Subset(0,\pi)$ keeps
$v_h$ away from zero in the chosen channel. Adding the three bounds of
\cref{eq:proto-lsz-bounds} gives \cref{eq:spec-structural}; the Adler
extension follows since the right-hand side is $O(\epsilon)$, and
$\arg(1+z)=\mathrm{Im}\,z+O(\abs{z}^2)$ with the nonzero lower bound on
$\abs{\bar k_*}$ gives \cref{eq:spec-phase-jet}.

\provenance{S-IDX-spec-struct-r2}{%
  \texttt{theory/soft-index-r2.md} §3, statement and proof at step (1.8);
  the limit-point register at (1.7)}{%
  no certificate for the hypothesis; the gates test the finite inputs only
  (\texttt{theory/checks/soft\_index\_r2\_check.py} gate SIDXR2-C5,
  arithmetic)}{%
  \texttt{theory/verdicts/soft-index-r2-critic.md}}

\subsection{Adding on-shell data: the matched value}

The structural law contains a shape and an open constant. To get a number one
must connect the protocol readout to something on shell. That connection is a
separate hypothesis, and it is displayed rather than assumed silently.

\begin{theorem}[Matched spectral law under the matching hypothesis]
\label{thm:soft-index-matched}
\statusSketch

Assume the hypotheses of \cref{thm:soft-index-structural}; a fully polarised
spin-$S$ vacuum, so that the tail density is $\rho=S$; a primitive one-magnon
hard leg, so that $\ell_h=1$; the class-membership antecedent
$(\beta_S)$ of \cref{thm:d24-val} --- existence of a source
$O\in\classSW(\rho)|_{\rho=S}$ with $M_1^O\neq0$ --- and the following separate
matching hypothesis:
\begin{equation}
  \textsc{(match-s)}:\qquad
  \Aproto_*(\epsilon)-\int[\,S_{\rm phys}(k,h)-1\,]\,d\mu_{*,\epsilon}(k,h)
  =o(\epsilon),
  \label{eq:match-s}
\end{equation}
a first-order identification of the fixed-time protocol readout with the
on-shell multiplier, containing no numerical slope of its own. Then every
actual limit point has phase slope
\begin{equation}
  \frac{\arg(1+\Aproto_*(\epsilon))}{\bar k_*(\epsilon)}
  \longrightarrow \frac{\sgn(v_h-v_s)}{S},
  \label{eq:matched-slope}
\end{equation}
and comparison with \cref{eq:spec-structural} supplies the conditional matched
value of the amputation constant,
\begin{equation}
  \aleg(S)=\frac{1}{2S}=\frac{1}{Z_\rho}.
  \label{eq:matched-value}
\end{equation}
More generally, any datum at the same density that shares this class constant
and satisfies \textsc{(proto-lsz)} obeys the structural slope
$\sgn(v_h-v_s)\ell_h/\rho$.
\end{theorem}

\begin{scope}[Where the number comes from, and what is not claimed]
\label{scope:soft-index-matched}
The only numerical supplier is the exact two-body expansion
$S_{\rm phys}(k,h)-1=i\sgn(v_h-v_s)k/S+O(k^2)$ of \cref{sec:two-magnon}, and
it enters strictly \emph{after} \textsc{(match-s)}, never as a clause of the
readout and never as a regularity condition. Hypothesis \textsc{(match-s)} is
exactly the open first-order bridge that the wave-operator results of
\cref{sec:wave-operators} decline to supply at this interface, and by
\cref{scope:regularity-no-value} the regularity package cannot substitute for
it; the theorem is therefore a complete conditional implication with no
instance in general.

The value \cref{eq:matched-value} at $S\in\{1/2,1,3/2,2\}$ is precisely the
existing conditional row of \cref{thm:d24-val}; at other half-integer spins it
is the same one-line consequence, but it belongs to this sketch rather than
being silently attributed to that row. Nothing here exhibits a member of
$\classSW(\rho)$, and nothing here licenses a bridge-free conclusion about the
amputation constant --- that would be the conjecture of \cref{conj:amp}. The
final displayed generalisation carries $\ell_h$ as the \emph{measured} Ward
register; no composite-charge formula $\ell_h=\abs{q}$ is asserted beyond the
primitive $\abs{q}=1$ case, since the two-body slope supplier is scoped to
unit charge. No leg-normalisation argument is used anywhere in the derivation
of \cref{eq:matched-value}: the exact charge-created/asymptotic leg conversion
supplies only $Z_\rho^{-1/2}$, and never $Z_\rho^{-1}$.
\end{scope}

\provenance{S-IDX-spec-r2}{%
  \texttt{theory/soft-index-r2.md} §3, steps (1.9).(2.1)--(2.5); the value
  interface and its fences at (1.11)}{%
  \texttt{theory/checks/soft\_index\_r2\_check.py} \texttt{-{}-red-s2-value}
  tests the finite matching arithmetic only; no gate tests \textsc{(match-s)}}{%
  \texttt{theory/verdicts/soft-index-r2-critic.md}}

\subsection{The separated-preparation breakthrough: matching becomes a theorem}

The results so far leave one hypothesis carrying all the weight. The advance
recorded here removes it --- not in general, but on an explicitly displayed
subclass of preparations, and there completely.

The idea is physical and simple. If the soft and hard packets are prepared far
apart, allowed to approach, collide once, and separate, and if the
separations and settling times are taken out in the right order at fixed soft
scale, then the fixed-time charge-created state can be compared \emph{in norm}
with a genuine incoming scattering channel, that channel can be pushed through
the two-body wave operators, and the readout can be evaluated on the outgoing
side. The comparison is a chain of Cook-type estimates; none of them needs to
be uniform in the soft scale, because the soft scale is held fixed until the
very end.

\begin{definition}[The separated-preparation subclass]
\label{def:separated-preparation}
Specialise the indices of \cref{def:protocol-readout} as follows. At fixed
$\epsilon>0$ take hard widths $\sigma_j\downarrow0$, initial separations
$R_j\to\infty$, and settling times $T_j$ with
$L_j:=d_\epsilon T_j-R_j\to\infty$, where $d_\epsilon$ is the velocity
separation of the two packets. With
$B_{M,\epsilon}$ the largest $\sup$-norm of the derivatives of
$S_{\rm phys}$ of total order at most $M$ on the soft support times $I$, and
$K_{M,\epsilon}:=C_{M,S,J}(1+B_{M,\epsilon})/d_\epsilon$, require for one
integer $M\ge8$
\begin{equation}
  K_{M,\epsilon}\;s_M(f_\epsilon)\;s_M(g_{\sigma_j})\,
  \bigl[(1+R_j)^{3-M}+(1+L_j)^{3-M}\bigr]\longrightarrow0,
  \label{eq:sep}
\end{equation}
where $s_M$ is the fixed-packet Schwartz seminorm evaluated on the
\emph{untranslated} envelopes, the unit-modulus translation phase being kept
inside the stationary phase rather than differentiated. The hard-column window
is the labelled rectangle formed by the soft support and $\supp g_{\sigma_j}$,
and zero-mass rows carry zero measure. After $(\sigma_j,R_j,T_j)$ are fixed,
choose a ring $N_j$ so large that $N_j\epsilon(c_2-c_1)>2\pi$,
$T_j\le c_{\rm rec}N_j/v_{\max}$, and the total norm error $\eta_j$ between
the periodised finite-ring kernels and their infinite-chain counterparts is at
most $1/j$.
\end{definition}

This class is nonempty and directly verifiable, not a renaming of the matching
hypothesis: compactly supported smooth packets satisfy
$s_M(g_\sigma)=O(\sigma^{-M-1/2})$, so taking, for instance,
$R_j=\lceil\sigma_j^{-2}\rceil$ and $T_j=2R_j/d_\epsilon$ makes $L_j=R_j$ and
the left-hand side of \cref{eq:sep} equal to $O_\epsilon(\sigma_j^{M-13/2})$,
which vanishes for $M\ge8$. Suitable rings exist by Plancherel sampling and
strong convergence of finite-range dynamics at fixed time, chosen diagonally.

\begin{theorem}[On the separated class, matching holds with zero remainder]
\label{thm:match-separated}
\statusProved

Let $H_S=-J\sum_x(\mathbf S_x\cdot\mathbf S_{x+1}-S^2)$ with $2S\in\N$, work on
the primitive regular two-magnon channel, and assume the separated-preparation
subclass of \cref{def:separated-preparation}. Then along every actual
row-measure limit subsequence the fixed-time charge-created finite-ring
readout converges and satisfies, for that same limiting row measure,
\begin{equation}
  \Aproto_*(\epsilon)=\int[\,S_{\rm phys}(k,h)-1\,]\;d\mu_{*,\epsilon}(k,h).
  \label{eq:match-separated}
\end{equation}
In particular the matching hypothesis \textsc{(match-s)} of
\cref{eq:match-s} holds on this subclass with \emph{zero} remainder, hence a
fortiori through $o(\epsilon)$.
\end{theorem}

\begin{scope}[The boundary of the positive result]
\label{scope:match-separated}
The theorem is a whole-datum scattering identity. It does \emph{not} identify
the finite-separation fixed-time vector with a Haag--Ruelle creator or with a
scattering vector; it does not require a soft-uniform Cook estimate; and it
proves no component decomposition. In particular it does not prove
\textsc{(proto-lsz)}: nothing above splits the datum into descendant,
orthogonal-current, direct-contact and two boundary-gradient pieces, exhibits
a microscopic member of $\classSW(\rho)$, or evaluates the amputation
constant. Compare \cref{eq:match-separated} with
\cref{eq:proto-lsz-split,eq:proto-lsz-bounds}: the latter is strictly stronger
content.

The result is primitive two-body only: no endpoint, equal-velocity,
composite-charge, bound-state-channel, many-body, or asymptotic-completeness
statement is asserted outside the two-magnon sector. The proof uses no scalar
replacement of the full-sector Gram operator, respecting the register erratum
of \cref{scope:soft-index-su2}.
\end{scope}

\paragraph{Idea of proof: a norm bridge.} Four movements. \emph{(i) Wave
operators.} Fourier transform in the centre coordinate decomposes the
two-magnon Hamiltonian into fibres over the total momentum; on each fibre the
relative motion is a constant-coefficient Jacobi recurrence away from contact,
with only the adjacent row $r=1$ and, for $S\ge1$, the single normalised
double-occupancy row $r=0$ modified. A finite-rank boundary perturbation of a
scalar Jacobi operator has isometric incoming and outgoing generalised
transforms with equal range, and their composition is multiplication by the
physical ratio, $(W_+^S)^*W_-^S=M_{S_{\rm phys}}$, with $S_{\rm phys}$ the
exact two-body multiplier of \cref{sec:two-magnon}. \emph{(ii) Cook
estimates.} The identification defect
$(H_SI_S-I_SH_0)U_0(t)F_R$ is supported only on the collision rows $r\le1$;
outside the two separated velocity cones repeated integration by parts gives
decay $(1+R+d_\epsilon\abs t)^{2-M}$, and integrating in $t$ gives the two
displayed bounds --- one comparing the fixed-time state with the incoming wave
operator, one comparing the outgoing channel with free propagation.
\emph{(iii) The bridge.} Equality of ranges gives
$W_-^SF_R=W_+^SM_{S_{\rm phys}}F_R$, and combining the two Cook bounds yields
a single norm estimate that starts at the fixed-time charge state and ends at
the freely propagated outgoing channel. No equality between those two states
was assumed anywhere. \emph{(iv) Through the readout.} The coordinate-kernel
map and the row-rectangle selection are contractions; the free-reference error
is the freely propagated initial collision rows plus a finite-ring sampling
error, both $o(1)$ under \cref{eq:sep}. Disjoint momentum supports make the
ideal labelled rectangles orthogonal of equal mass $1/2$, which bounds the
denominators away from zero, so the row measures converge in total variation
and the normalised pairing converges. In the ideal rectangle the free
propagation phase and the translation phase occur on both sides of the pairing
and cancel, leaving exactly $\int S_{\rm phys}\,d\mu$. Since the difference in
\textsc{(match-s)} is identically zero at each fixed $\epsilon$ after the
outer limit, it is trivially $o(\epsilon)$ when $\epsilon\downarrow0$ is
finally taken --- which is precisely why the fixed-packet divergence of the
Cook constants does not obstruct this subclass.

\provenance{S-IDX-MATCH-HS-SEP}{%
  \texttt{theory/proto-lsz-match.md} steps (1.1)--(1.6): the subclass (1.2),
  wave operators (1.3), the fixed-time-to-incoming-channel bound (1.4),
  passage through the readout (1.5), the theorem (1.6); honest boundary at
  (1.11)}{%
  \texttt{theory/checks/proto\_lsz\_match\_check.py} gates PMLM-C1--C5; finite
  corroboration only, with no wave-operator certificate}{%
  \texttt{theory/verdicts/blitz-proto-lsz-r1.md}}

\begin{theorem}[The value on the separated class]
\label{thm:d29-value-separated}
\statusProved

Under the separated-preparation subclass of
\cref{def:separated-preparation}, along every actual row-measure limit
subsequence the primitive fixed-time scalar datum obeys
\begin{equation}
  \Aproto_*(\epsilon)=\frac{i\,\sgn(v_h-v_s)\,\bar k_*(\epsilon)}{S}
  +O(\epsilon^2),
  \qquad
  \frac{\arg(1+\Aproto_*(\epsilon))}{\bar k_*(\epsilon)}
  \longrightarrow\frac{\sgn(v_h-v_s)}{S}.
  \label{eq:d29-value-sep}
\end{equation}
\end{theorem}

\begin{scope}[This row is about the readout, not about the class constant]
\label{scope:d29-value-separated}
The value enters only after \cref{thm:match-separated} has already identified
the readout with the multiplier, and then solely through the exact two-body
slope of \cref{sec:two-magnon}. The row makes no claim about the amputation
constant $\aleg$, about membership in the no-contact class of
\cref{def:no-contact-class}, about hypothesis \textsc{(proto-lsz)}, or about
endpoints, equal velocities, composite charges, or any other protocol class.
In particular, no value of $\aleg$ follows from this corollary alone; if
\textsc{(proto-lsz)}, class membership, and the other antecedents of
\cref{thm:soft-index-matched} were separately supplied, the comparison there
would give the conditional value $\aleg(S)=1/(2S)$ --- this theorem neither
supplies nor replaces those antecedents.
\end{scope}

\paragraph{Idea of proof.} Insert
$S_{\rm phys}(k,h)-1=i\sgn(v_h-v_s)k/S+O(k^2)$ into
\cref{eq:match-separated} and integrate against the same measure; the
scale-tied support $\abs k\le c_2\epsilon$ makes the integrated remainder
$O(\epsilon^2)$. Then $\arg(1+z)=\mathrm{Im}\,z+O(\abs z^2)$ together with the
one-sided lower bound $\abs{\bar k_*}\ge c_1\epsilon$ gives the slope.

\provenance{S-IDX-D29-value-HS-SEP}{%
  \texttt{theory/proto-lsz-match.md} step (1.10) --- the only place in that
  shard where the number $1/S$ is invoked}{%
  \texttt{theory/checks/proto\_lsz\_match\_check.py} gate PMLM-C4 plus the
  frozen partial gate C5}{%
  \texttt{theory/verdicts/blitz-proto-lsz-r1.md}}

\subsection{An exact LSZ-shaped scalar factorisation}

The matched identity can be rewritten so that it has the \emph{shape} the
universality analysis of \cref{sec:universality} predicts --- an external flux
factor times a Ward residue, integrated against the row measure --- with a
reduced multiplier that extends continuously to zero soft momentum. This is
worth isolating, because it shows that the LSZ shape itself is available on
the separated class without any of the component hypotheses.

\begin{theorem}[Exact scalar factorisation of every separated limit point]
\label{thm:proto-scalar}
\statusProved

On the fully polarised bilinear ferromagnet $H_S$ with $2S\in\N$, and on the
separated-preparation subclass of \cref{def:separated-preparation}, define for
$k\neq0$ the \emph{scalar scattering witness}
\begin{equation}
  L_S^{\rm sc}(k,h):=\frac{S_{\rm phys}(k,h)-1}{(e^{ik}-1)\,2iv_h}.
  \label{eq:scalar-witness}
\end{equation}
Then $L_S^{\rm sc}$ has a unique uniformly $C^1$ extension to $k=0$ on the
compact channel, with
\begin{equation}
  L_S^{\rm sc}(0,h)=\frac{-i\,\sgn(v_h-v_s)}{2S\,v_h},
  \label{eq:scalar-witness-value}
\end{equation}
and every actual row-measure limit point has the exact factorisation
\begin{equation}
  \Aproto_*(\epsilon)=
  \int(e^{ik}-1)\,L_S^{\rm sc}(k,h)\,[\,2iv_h\,]\;d\mu_{*,\epsilon}(k,h),
  \qquad
  \Aproto_*(\epsilon)=\frac{i\,\sgn(v_h-v_s)\,\bar k_*(\epsilon)}{S}
  +O(\epsilon^2).
  \label{eq:scalar-factorisation}
\end{equation}
\end{theorem}

\begin{scope}[Same profile is not same object]
\label{scope:proto-scalar}
Equation \cref{eq:scalar-witness-value} has the same functional profile as the
no-contact class clause \cref{def:no-contact-class}(3b) with scalar
coefficient $1/(2S)$. That observation does \emph{not} identify
$L_S^{\rm sc}$ with the descendant quotient $L$ of that clause, and it does
not prove a member of $\classSW(\rho)|_{\rho=S}$: the class's $L$ is the
quotient of an independently identified descendant external-leg component
inside an exhaustive decomposition, and is defined only when the class is
nonempty. The superscript is a deliberate warning label. Nor does the theorem
prove the orthogonal-current, direct-contact, or window bounds of
\textsc{(proto-lsz)}, which remains open. No step obtains $1/(2S)$ from the
charge-created/asymptotic-leg conversion, which supplies only
$Z_\rho^{-1/2}$.
\end{scope}

\begin{scope}[The exact remaining lemma \textsc{(comp-hs)}]
\label{scope:comp-hs}
Aggregate matching cannot determine component estimates, and the obstruction
is elementary: for any nonzero bounded smooth $b$ on $I$ and
$B(k,h)=k\,b(h)$, the replacement
$E_{\rm desc}\mapsto E_{\rm desc}+B$, $E_\perp\mapsto E_\perp-B$ leaves every
whole-datum equality unchanged while moving both components. What is therefore
needed, and all that is needed, is one model lemma.

\textsc{(comp-hs)}: on the separated class, using at finite volume the exact
full-sector split $J^-_k=P_{\lambda,N}J^-_k+(1-P_{\lambda,N})J^-_k$ and the
exact Duhamel contact defect --- with no scalar replacement of the Gram
operator --- the finite readout functional admits four independently defined
terms, from the descendant-range current projection, the complementary
projection, the direct source/contact functional, and the two window
gradients, such that on the ordered outer limit: \emph{(1)} the descendant
quotient converges in the packet norm to a uniformly $C^1$ multiplier
$L_S^{\rm desc}$ with $L_S^{\rm desc}(0,h)=L_S^{\rm sc}(0,h)$, equivalently
$L_S^{\rm desc}-L_S^{\rm sc}=O(k)$ uniformly; \emph{(2)} the
complementary-current and direct-contact terms are each $O(\epsilon^2)$
uniformly on $I$; \emph{(3)} the two window-gradient terms sum to
$o(\epsilon)$ under \cref{eq:sep}; \emph{(4)} the limiting descendant
construction exhibits a nonzero microscopic member of
$\classSW(\rho)|_{\rho=S}$.

Together with \cref{thm:proto-scalar} this would prove full
\textsc{(proto-lsz)} on the separated class without changing any definition.
No existing proved claim supplies it: the finite Ward theorem gives no
volume-uniform reduced-channel regularity, and the matching theorem controls
only the sum after the outer limit. A viable attack is a two-body fibre
limiting-absorption or Feshbach estimate --- insert the full-sector descendant
projector before taking the Jacobi boundary resolvent, prove a uniform
reduced-resolvent bound on the compact hard window, and push the four finite
functionals through the already proved norm bridge --- which tests items (1)--(3)
directly rather than inferring them from their sum.
\end{scope}

\paragraph{Idea of proof.} The multiplier $S_{\rm phys}$ is analytic in $k$
and smooth in $h$ near the compact rectangle with $S_{\rm phys}(0,h)=1$ and
$\partial_kS_{\rm phys}(0,h)=i\sgn(v_h-v_s)/S$. The analytic removable-factor
theorem therefore gives $S_{\rm phys}-1=(e^{ik}-1)G_S(k,h)$ with $G_S$
uniformly $C^1$ and $G_S(0,h)=\sgn(v_h-v_s)/S$, since
$\partial_k(e^{ik}-1)|_0=i$; dividing by $2iv_h$, which is bounded away from
zero on $I$, preserves regularity and gives
\cref{eq:scalar-witness-value}. Substituting the resulting pointwise identity
$S_{\rm phys}-1=(e^{ik}-1)L_S^{\rm sc}\,[2iv_h]$ --- valid at $k=0$ by
continuity --- into the matched identity \cref{eq:match-separated} gives
\cref{eq:scalar-factorisation}. The number $1/(2S)$ entered only through the
$k$-derivative of the physical multiplier, after its branch was fixed.

\provenance{S-IDX-PROTO-SCALAR-HS-SEP}{%
  \texttt{theory/proto-lsz-scalar.md} steps (1.1)--(1.7): removable quotient
  (1.2), exact factorisation (1.3), supplier separation (1.4), the
  aggregate/component obstruction (1.5), the lemma \textsc{(comp-hs)} (1.6)}{%
  \texttt{theory/checks/proto\_lsz\_check.py} gates PLSZ-C1--C4 (corroboration
  of the scalar algebra only, not of regularity)}{%
  \texttt{theory/verdicts/blitz-proto-lsz-r1.md} (0 fatal, 0 major)}

\subsection{The Haag--Ruelle value instance}

Finally, there is a second and logically independent place where the number
$1/S$ appears: not in the fixed-time protocol readout at all, but in the
constructed scattering channel of \cref{sec:wave-operators}. It is recorded
here to keep the two clearly apart.

\begin{theorem}[Value in the constructed Haag--Ruelle channel]
\label{thm:hr-value}
\statusSketch

Under the exact-band register on the ferromagnetic model, the fixed-packet
constructed Haag--Ruelle channel of \cref{sec:wave-operators} has multiplier
$S_{\rm phys}$, and its scale-tied packet average
\begin{equation}
  \Aproto_{\rm HR}(\epsilon):=
  \int[\,S_{\rm phys}(\epsilon u,h)-1\,]\,d\mu_f(u,h)
  \label{eq:hr-average}
\end{equation}
has phase slope $\sgn(v_h-v_s)/S$, with
$\bar k=\epsilon\int u\,d\mu_f$.
\end{theorem}

\begin{scope}[An on-shell instance, explicitly not a protocol instance]
\label{scope:hr-value}
Equation \cref{eq:hr-average} is a value instance for a \emph{constructed}
channel. It is not the fixed-time charge-created protocol datum of
\cref{def:protocol-readout}: the Cook and Gram constants of the wave-operator
construction diverge along the rescaled soft packet, and creator-choice
independence does not apply to the fixed-time insertion
$Q[f_\epsilon]\psi_g$. It is therefore no basis for upgrading either the
structural or the matched spectral row. Consequently there is at present no
proved instance of \cref{thm:soft-index-structural} or
\cref{thm:soft-index-matched} for the general protocol, no proved microscopic
member of $\classSW(\rho)$, and no permission to use the amputation conjecture
of \cref{conj:amp} as a value supplier.
\end{scope}

\provenance{S-IDX-HR-value-r2}{%
  \texttt{theory/soft-index-r2.md} §4, instances and non-instances at step
  (1.10); interfaces and fences at (1.11)}{---}{%
  \texttt{theory/verdicts/soft-index-r2-critic.md}}

\subsection{Reading of the section}

The symmetry algebra alone delivers a genuine index: exact, finite, integral
where an index should be integral, and valid for every compact group one root
at a time. That is \cref{thm:soft-index-su2,thm:soft-index-compact-G,%
thm:sector-label}, and none of it is conditional.

Passing to the soft limit costs one hypothesis --- an exhaustive component
decomposition of the limiting readout --- and passing from a shape to a number
costs a second --- the identification of that readout with the on-shell
multiplier. On the separated-preparation class the second cost is paid in
full: matching is a theorem there, with zero remainder, and the value and the
LSZ shape follow. The first cost is not yet paid anywhere, and it is now
isolated to the single lemma \textsc{(comp-hs)} of \cref{scope:comp-hs}. That
lemma is the one remaining gap of the one-dimensional soft-index programme.
